\section{Related Literature}\label{sec:literature}

The paper connects three literatures. The contribution is the interaction among them, not a claim to introduce any one component.

First, a long literature studies strategic local public inputs and mobile firms. \citet{WalzWellisch1996} analyze strategic public provision of local inputs with endogenous oligopolistic location choice. \citet{MaurerWalz2000} study two regional governments competing for two mobile oligopolistic firms and show that regional competition generally produces suboptimal infrastructure provision. Related work shows that mobility and market integration can change the direction of local-public-good distortions \citep{Matsumoto1998,Ihara2008}. More recently, \citet{MengDingWang2024} document and model a negative relationship between local-government competition and industrial resilience operating through market segmentation and industrial-structure distortions. Our plant-attraction motive is deliberately close to this literature. The difference is that public protection also changes firms' privately chosen geographic redundancy, so the planner's marginal valuation of the same local input becomes endogenous to the industrial resilience response.

Second, environmental and regional models make firm location and market structure endogenous to policy. \citet{MarkusenMoreyOlewiler1993} show that environmental policy can induce discrete changes in plant location and market structure, and \citet{MarkusenMoreyOlewiler1995} study noncooperative regional environmental policy when plant locations are endogenous. Closest on strategic environmental policy and industrial allocation, \citet{MartinHerranEtAl2026} analyze two regions in a transboundary-pollution dynamic game in which environmental policy affects the long-run spatial distribution of firms. Their analysis emphasizes pollution policy, transport costs, and agglomeration/dispersion forces. Our mechanism instead treats public protection against disruption as the policy instrument and makes firms' costly geographic backup readiness endogenous; the resulting welfare channel runs through disaster-state production availability and market structure. \citet{OtazawaEtAl2016} incorporate natural-disaster risk into an economic-geography environment and show that market equilibrium may display excessive industrial agglomeration and vulnerability. Especially close to one part of our mechanism, \citet{GramesEtAl2019} study a firm's joint flood-risk location and protection decisions and show that existing flood defense can induce the firm to locate closer to water and invest less in private flood defense. We therefore do not claim novelty for policy-induced location change, adaptation crowd-out, or the joint occurrence of those two responses for a single firm.

Third, recent work studies private adaptation and resilience. \citet{Eisenack2014} analyzes private adaptation with endogenous market structure. \citet{CapponiDuStiglitz2024} show that market power and market incompleteness can produce insufficient supply-network resilience. \citet{CastroVincenziEtAl2024} study firms' geographic diversification of sourcing under climate risk. Most directly on public/private substitution, a 2026 working paper by Zhao, Yang, and Zhang (SSRN 6996377) reports evidence that public flood protection can crowd out private defensive investment and increase tail-risk losses. Accordingly, public/private resilience substitution is not our novelty claim.

The paper's contribution is narrower: a full-game policy-ranking result produced by combining protection-induced private geographic-redundancy adjustment, mobile oligopolistic plant location, and decentralized competition for plants. In the fixed-redundancy benchmark, the model resembles local-public-input competition. In the fixed-location benchmark, public/private resilience substitution remains but the plant-attraction motive disappears. In the full game, Proposition~\ref{prop:attraction-threshold} shows transparently how the local plant-attraction term can overturn a negative national marginal value of additional protection, while Theorem~\ref{thm:main} establishes a global policy-ranking reversal on a nonempty open set. This policy-ranking reversal, rather than any single ingredient, is the paper's novelty claim.
